
\documentclass{article}
\usepackage{colortbl}
\usepackage{makecell}
\usepackage{multirow}
\usepackage{supertabular}

\begin{document}

\newcounter{utterance}

\centering \large Interaction Transcript for game `cladder', experiment `full\_v1.5\_default', episode 1468 with Qwen3.6{-}27B.
\vspace{24pt}

{ \footnotesize  \setcounter{utterance}{1}
\setlength{\tabcolsep}{0pt}
\begin{supertabular}{c@{$\;$}|p{.15\linewidth}@{}p{.15\linewidth}p{.15\linewidth}p{.15\linewidth}p{.15\linewidth}p{.15\linewidth}}
    \# & \multicolumn{2}{c}{Player} && \multicolumn{2}{c}{Game Master} \\
    \hline

    \theutterance \stepcounter{utterance}  
    & & & \multicolumn{4}{p{0.6\linewidth}}{
        \cellcolor[rgb]{0.9,0.9,0.9}{
            \makecell[{{p{\linewidth}}}]{
                \texttt{\tiny{[P1$\langle$GM]}}
                \texttt{You are an expert in causal inference. The following question is not a typical commonsense query, but rather a meticulously designed question created by a professor specializing in causal inference, intended to assess the students' mastery of the course content.} \\
\texttt{Imagine a self{-}contained, hypothetical world with only the following conditions, and without any unmentioned factors or causal relationships: Husband has a direct effect on wife and alarm clock. Wife has a direct effect on alarm clock. The overall probability of alarm set by husband is 8\%. For husbands that don't set the alarm, the probability of ringing alarm is 75\%. For husbands that set the alarm, the probability of ringing alarm is 16\%. Is ringing alarm more likely than silent alarm overall?} \\
\texttt{Start your answer with 'Yes' or 'No', followed by additional reasoning or evidence to support your explanation.} \\
            }
        }
    }
     \\ \\

    \theutterance \stepcounter{utterance}  
    & \multicolumn{4}{p{0.6\linewidth}}{
        \cellcolor[rgb]{0.561,0.737,0.561}{
            \makecell[{{p{\linewidth}}}]{
                \texttt{\tiny{[P1$\rangle$GM]}}
                \texttt{No} \\
\\ 
\texttt{To determine whether the ringing alarm is more likely than the silent alarm overall, we must calculate the total probability of the alarm ringing, denoted as \$P(R=1)\$, where \$R=1\$ represents the alarm ringing and \$R=0\$ represents the alarm being silent.} \\
\\ 
\texttt{**1. Identify the Variables and Probabilities:**} \\
\texttt{*   Let \$H\$ be the event that the husband sets the alarm.} \\
\texttt{*   Let \$R\$ be the event that the alarm rings.} \\
\texttt{*   We are given the marginal probability of the husband setting the alarm: \$P(H=1) = 0.08\$.} \\
\texttt{*   Consequently, the probability that the husband does *not* set the alarm is: \$P(H=0) = 1 {-} 0.08 = 0.92\$.} \\
\texttt{*   We are given the conditional probability of the alarm ringing given the husband did *not* set it: \$P(R=1 | H=0) = 0.75\$.} \\
\texttt{*   We are given the conditional probability of the alarm ringing given the husband *did* set it: \$P(R=1 | H=1) = 0.16\$.} \\
\\ 
\texttt{*Note: The problem states that "Wife has a direct effect on alarm clock," but it does not provide any specific probabilities or conditional distributions involving the wife's actions. In causal inference problems of this type, if specific conditional probabilities} \\
            }
        }
    }
    & & \\ \\

    \theutterance \stepcounter{utterance}  
    & & & \multicolumn{2}{p{0.3\linewidth}}{
        \cellcolor[rgb]{0.9,0.9,0.9}{
            \makecell[{{p{\linewidth}}}]{
                \texttt{\tiny{[GM$|$GM]}}
                \texttt{No} \\
            }
        }
    }
    & & \\ \\

    \theutterance \stepcounter{utterance}  
    & & & \multicolumn{2}{p{0.3\linewidth}}{
        \cellcolor[rgb]{0.9,0.9,0.9}{
            \makecell[{{p{\linewidth}}}]{
                \texttt{\tiny{[GM$|$GM]}}
                \texttt{yes} \\
            }
        }
    }
    & & \\ \\

    \theutterance \stepcounter{utterance}  
    & & & \multicolumn{2}{p{0.3\linewidth}}{
        \cellcolor[rgb]{0.9,0.9,0.9}{
            \makecell[{{p{\linewidth}}}]{
                \texttt{\tiny{[GM$|$GM]}}
                \texttt{game\_result = LOSE} \\
            }
        }
    }
    & & \\ \\

\end{supertabular}
}

\end{document}
